\documentclass[10pt,a4paper]{article}

% --- Packages ---
\usepackage[utf8]{inputenc}
\usepackage{mathptmx} % Times New Roman
\usepackage[T1]{fontenc}
\usepackage{amsmath}
\usepackage{graphicx}
\usepackage{caption}
\usepackage{float}
\usepackage{hyperref}
\usepackage{xcolor}
\usepackage{authblk}
\usepackage{booktabs}
\usepackage{tabularx}
\usepackage{titlesec}
\usepackage{tcolorbox}
\usepackage{multirow}
\usepackage{fancyhdr}
\usepackage{geometry}
\geometry{a4paper, top=1.3cm, bottom=1.3cm, left=2cm, right=2cm, includeheadfoot}

% --- Dark Blue Running Header & Footer (ITEGAM-JETIA style) ---
\pagestyle{fancy}
\fancyhf{}
\definecolor{navyblue}{HTML}{001a57}
\setlength{\headheight}{28pt}
\setlength{\headsep}{8pt}
\setlength{\footskip}{32pt}
\fancyhead[C]{%
  \colorbox{navyblue}{\makebox[\textwidth][c]{\rule[-4pt]{0pt}{18pt}%
    \textcolor{white}{\normalsize\textbf{ITEGAM-JETIA, Manaus, v.XX n.XX, p.XX-XX. Mon./Mon., 20XX.}}}}
}
\fancyfoot[C]{%
  \colorbox{navyblue}{\makebox[\textwidth][c]{\rule[-4pt]{0pt}{18pt}%
    \textcolor{white}{\normalsize\textbf{Page \thepage}}}}
}
\renewcommand{\headrulewidth}{0pt}
\renewcommand{\footrulewidth}{0pt}

% --- Hyperref Setup ---
\hypersetup{
    colorlinks=true,
    linkcolor=blue,
    filecolor=magenta,
    urlcolor=blue,
    citecolor=blue,
}

% --- Heading Formatting (ITEGAM-JETIA style) ---
\titleformat{\section}{\normalfont\normalsize\bfseries\color{navyblue}\uppercase}{\thesection.}{1em}{}
\titleformat{\subsection}{\normalfont\normalsize\bfseries\color{navyblue}}{\thesubsection}{1em}{}
\titleformat{\subsubsection}{\normalfont\normalsize\color{navyblue}}{\thesubsubsection}{1em}{}

\begin{document}

% --- Header / Journal Info ---
\begin{center}
    {\Large \textbf{Journal of Engineering and Technology for Industrial Applications}} \\
    \vspace{0.2cm}
    {\huge \textbf{ITEGAM-JETIA}} \\
    \vspace{0.2cm}
    \colorbox{yellow}{\textbf{Manaus, v.XX n.XX, p. XX-XX. Mon./Mon., 20XX.}} \\
    \colorbox{yellow}{\textbf{DOI: https://doi.org/10.5935/jetia.XXXXX.XXX}}
\end{center}

\vspace{0.3cm}
\noindent\colorbox{navyblue}{\makebox[\textwidth][l]{\textcolor{white}{\textbf{RESEARCH ARTICLE \hfill OPEN ACCESS}}}}
\vspace{0.5cm}

% --- Title Block ---
\vspace{0.4cm}
\noindent{\color{navyblue}\rule{\textwidth}{1.5pt}}
\vspace{0.3cm}

\begin{center}
    {\Large \textbf{\color{navyblue}KRISHI-DRISHTI: SATELLITE-DRIVEN CARBON CREDIT VERIFICATION FOR SMALLHOLDER FARMERS USING SENTINEL-2 REMOTE SENSING, FOUR-PILLAR dMRV FRAMEWORK, AND CRYPTOGRAPHIC HARVEST TOKENIZATION}}
\end{center}

\vspace{0.2cm}
\noindent{\color{navyblue}\rule{\textwidth}{2.5pt}}
\vspace{0.4cm}

\begin{center}
    {\normalsize \textbf{Dr. Rama Krishna K*$^1$, Dr. Sheethal Aji Mani$^2$, Bhagesh Biradar$^3$, Manjunath Khot$^4$, Jayanth R$^5$ and Pruthviraja S$^6$}} \\
    \vspace{0.4cm}
    {\small $^{1,2,3,4,5,6}$Department of CSE Data Science, Impact College of Engineering and Applied Sciences, Visvesvaraya Technological University, Bangalore, Karnataka, India.} \\
    \vspace{0.3cm}
    {\footnotesize
    $^1$\url{https://orcid.org/0000-0000-0000-0001},
    $^2$\url{https://orcid.org/0000-0000-0000-0002},
    $^3$\url{https://orcid.org/0000-0000-0000-0003}, \\
    $^4$\url{https://orcid.org/0000-0000-0000-0004},
    $^5$\url{https://orcid.org/0000-0000-0000-0005},
    $^6$\url{https://orcid.org/0000-0000-0000-0006}} \\
    \vspace{0.2cm}
    {\small Email: *ramakrishna@example.com, sheethal@example.com, bhagesh@example.com, manjunath@example.com, jayanth@example.com, pruthviraja@example.com}
\end{center}

\vspace{0.4cm}

% --- Article Info & Abstract ---
\noindent\begin{tabularx}{\textwidth}{@{}p{0.3\textwidth} X@{}}
\toprule
\textbf{ARTICLE INFO} & \textbf{ABSTRACT} \\
\midrule
\textbf{Article History} \newline
Received: Month 01th, 20XX \newline
Reviewed: Month 05th, 20XX \newline
Accepted: Month 15th, 20XX \newline
Published: Month 30th, 20XX
&
More than 86\% of India's 600 million farmers operate on plots smaller than two hectares. These smallholder farmers are largely excluded from carbon credit markets because verification costs can surpass USD \$3,000 per farm. We present Krishi-Drishti (``Farmer's Vision''), a mobile-first platform built around a satellite-driven Digital Measurement, Reporting, and Verification (dMRV) system. It allows smallholders to enroll in and monetize soil carbon credits without relying heavily on expensive lab tests or external auditors. By combining Sentinel-2 Level-2A multispectral imagery from Google Earth Engine (GEE) with a stacking ensemble machine learning approach, the system estimates Soil Organic Carbon (SOC) across fragmented plots. The credit calculation aligns with Verra VM0042, India's CCTS/BEE, and Gold Standard LUF AGR FM, enforcing a Four-Pillar MRV evidence framework: (1) land boundary and practice additionality verification; (2) IPCC Tier 2 physical soil sampling limits; (3) continuous satellite SOC monitoring; and (4) cryptographic verification via a HarvestToken ledger. Validated on 50 physical soil plots, the satellite SOC ensemble achieved an R$^2$ of 0.89 and an RMSE of 1.52 g/kg. Furthermore, a district-level additionality scoring engine showed 91.7\% classification accuracy. The platform also offers AI-based crop disease diagnosis (MobileNetV2, 88.1\% accuracy), bioacoustic pest detection, and a peer-to-peer marketplace. Usability testing with five farmers resulted in an 84.0/100 System Usability Scale (SUS) score. Overall, the system facilitates an estimated average income supplement of USD \$52--\$132 per farm annually, demonstrating how satellite-verified carbon sequestration can support rural livelihoods. \\
\addlinespace
\textbf{Keywords:} \newline
Satellite Carbon Verification, \newline
Soil Organic Carbon, \newline
Verra VM0042 / CCTS, \newline
Digital MRV (dMRV), \newline
HarvestToken, \newline
Sentinel-2 Remote Sensing.
& \\
\bottomrule
\end{tabularx}

\vspace{0.3cm}
\noindent\fbox{%
    \parbox{\dimexpr\textwidth-2\fboxsep-2\fboxrule\relax}{%
        Copyright \copyright 2026 by authors and Galileo Institute of Technology and Education of the Amazon (ITEGAM). This work is licensed under the Creative Commons Attribution International License (CC BY 4.0).
    }%
}
\vspace{0.5cm}

% ═══════════════════════════════════════════════════════════════════════════════
\section{INTRODUCTION}
% ═══════════════════════════════════════════════════════════════════════════════

Agriculture forms the core of India's economy and supports around 600 million people. However, the sector is under significant pressure: global food production needs to grow by over 70\% by 2050 to feed 10 billion people [1], yet agriculture already accounts for nearly 11\% of global greenhouse gas (GHG) emissions [3]. To meet both production and climate goals, managing Soil Organic Carbon (SOC) is essential. Regenerative practices like no-till farming, cover cropping, and agroforestry not only increase soil carbon sequestration and restore soil health, but also help generate tradeable carbon credits and boost overall crop yields.

Despite these benefits, smallholder farmers are often excluded from carbon markets due to technological and financial barriers. Conventional carbon credit verification requires physical soil laboratory testing, which can cost between USD \$1,500 and \$3,000 per farm. Combined with on-site audits and complex documentation processes, verification can take months [4]. Consequently, the global voluntary carbon market, valued at USD \$2 billion in 2023 [14], is largely out of reach for India's 600 million farmers, limiting their ability to participate in and benefit from climate mitigation efforts.

India's Carbon Credit Trading Scheme (CCTS), introduced under the Energy Conservation (Amendment) Act, 2022 and administered by the Bureau of Energy Efficiency (BEE), establishes a domestic regulatory framework for carbon credit trading [3]. Meanwhile, the Verra VM0042 Improved Agricultural Land Management (IALM) methodology provides the global technical standard for quantifying SOC sequestration from regenerative land management practices. Both frameworks require rigorous Measurement, Reporting, and Verification (MRV) — the exact bottleneck that excludes smallholders.

The availability of free, high-resolution multispectral satellite data—such as imagery from the European Space Agency's Sentinel-2 constellation—alongside cloud-based geospatial tools like Google Earth Engine (GEE), offers a scalable way to reduce verification costs. Research shows that Sentinel-2 spectral indices, notably the Bare Soil Index (BSI), red-edge bands, and Short Wave Infrared (SWIR) channels, closely correlate with SOC content in topsoils [6, 10, 12]. Using these satellite-derived SOC proxies instead of physical lab tests makes it technically possible to deliver accurate, continuous carbon credit assessments to smallholder farmers directly through their smartphones.

This paper introduces Krishi-Drishti (``Farmer's Vision''), a mobile-first platform focused on a satellite-driven Digital MRV (dMRV) system for verifying agricultural carbon credits. The core component, the Carbon Vault, uses a Four-Pillar MRV evidence framework to automate farm enrollment, verification, and credit issuance in alignment with Verra VM0042, CCTS/BEE, and Gold Standard LUF AGR FM. To lower costs and secure data integrity, issued credits are recorded as HarvestTokens using SHA-256 hash chaining. This ensures a tamper-evident ledger that meets the audit requirements of international carbon buyers and CBAM-compliance bodies. Furthermore, Krishi-Drishti incorporates AI-based crop disease diagnosis, bioacoustic pest detection, and a peer-to-peer produce marketplace to maintain farmer engagement.

The main contributions of this paper are as follows:
\begin{itemize}
    \item \textbf{[Primary]} A satellite-driven Four-Pillar Digital MRV (dMRV) system aligned with Verra VM0042, India's CCTS/BEE, and Gold Standard LUF AGR FM, allowing smallholder farmers to enroll and verify soil carbon credits via a mobile device without extensive physical audits.
    \item \textbf{[Primary]} A Sentinel-2-based ensemble SOC prediction pipeline (Random Forest, XGBoost, and CatBoost) achieving R$^2$ = 0.89 and RMSE = 1.52 g/kg on 50 geo-referenced Indian farmland plots.
    \item \textbf{[Primary]} A cryptographic HarvestToken architecture that uses SHA-256 hash chaining and daily Merkle Root Anchors for an immutable, CBAM-compliant carbon credit ledger without needing a live blockchain node.
    \item \textbf{[Primary]} A district-level additionality scoring engine, based on ICAR and NABARD adoption rate data, that classifies regenerative practices with 91.7\% accuracy under Verra/CCTS thresholds.
    \item \textbf{[Supporting]} The integration of a MobileNetV2 CNN (88.1\% accuracy) with a Google Gemini 2.5 Flash LLM for multi-lingual crop disease diagnosis, tied to carbon program recommendations.
    \item \textbf{[Supporting]} Additional tools, including a bioacoustic pest detection module and a peer-to-peer Mandi marketplace, to support overall rural livelihoods.
\end{itemize}

\begin{figure}[H]
    \centering
    \fbox{\includegraphics[width=0.95\textwidth]{fig_system_overview.jpg}}
    \caption{Krishi-Drishti Platform Architecture: The Carbon Vault (centre, highlighted) acts as the primary module integrating Sentinel-2 satellite data, the Farmer Mobile App, Admin Dashboard, MobileNetV2 AI, Mandi Marketplace, and Bioacoustic Detector. \newline \centering Source: Authors, (2026).}
\end{figure}

% ═══════════════════════════════════════════════════════════════════════════════
\section{THEORETICAL REFERENCE}
% ═══════════════════════════════════════════════════════════════════════════════

\subsection{Carbon Credit Markets and Agricultural MRV Frameworks}
The Verra VM0042 Improved Agricultural Land Management (IALM) methodology establishes the international standard for quantifying greenhouse gas emission reductions and SOC sequestration from improved land management. It requires three key evidential elements: a documented baseline emission scenario, a verifiable increase in soil carbon from an ``additional'' regenerative practice, and correction deductions for leakage ($LE_y$) and permanence risk (buffer pool). Gokul et al. (2026) established that India's CCTS — the domestic compliance framework under the BEE — offers a scalable alternative pathway for agricultural carbon offsetting within India, though high verification costs remain a barrier for smallholders [3]. The voluntary carbon market, governed by Verra VCS and Gold Standard, currently prices soil carbon removal credits at USD \$15--\$25 per tonne CO$_2$e, rising to USD \$50+ for higher-integrity nature-based solutions [14].

Digital MRV (dMRV) platforms provide a scalable solution to the smallholder inclusion problem. Organizations such as Boomitra and Grow Indigo have shown that combining satellite-derived SOC proxies with minimal physical soil sampling can reduce per-farm verification costs substantially. An ongoing challenge in these systems is the mathematical validation of the satellite-to-carbon mapping. The IPCC Tier 2 SOC stock formula provides the physical ground truth against which satellite models must be calibrated:
\begin{equation}
    SOC_{stock}\,(tC/ha) = \frac{SOC_{\%}}{100} \times BD\,(g/cm^3) \times Depth\,(m) \times 10
\end{equation}
Krishi-Drishti incorporates this formula directly into its verification engine, with administrative cross-validation of submitted laboratory certificates to deter data fraud.

\subsection{Satellite Remote Sensing for Soil Organic Carbon Prediction}
Accurate satellite-based SOC prediction is foundational to dMRV systems. Khaddi et al. (2026) indicate that combining Sentinel-2 multispectral data with Machine Learning models — particularly ensemble methods — often outperforms single-algorithm approaches for spatial SOC mapping [10]. Xiao et al. (2024) found that adding environmental covariates (terrain, climate, soil texture) alongside Landsat spectral data in an SVM model resulted in an R$^2$ of 0.9590, highlighting the importance of multi-feature inputs for modeling accuracy [11].

Iswarya \& Tamizhazhagan (2026) applied Quantile Random Forests (QRF) to Sentinel-2 indices across Tamil Nadu farmlands, finding strong predictive accuracy using red-edge (B5, B6, B7) and SWIR (B11, B12) bands [12]. Vaudour et al. (2022) suggested temporal mosaicking and bare-soil NDVI thresholding as essential preprocessing steps [6]. Triantakonstantis \& Karakostas (2025) further noted that stacking ensemble frameworks tend to yield the best SOC prediction results with Sentinel-2 imagery [15], which guided our decision to use a stacking ensemble architecture.

\subsection{Additionality Assessment in Carbon Credit Systems}
Additionality is a core requirement for carbon credit certification. A practice is considered ``additional'' only if it would not have occurred without the incentive provided by the carbon credit. Without additionality, credits do not represent genuine emission reductions [14]. Both Verra VM0042 and the CCTS/BEE framework dictate that a practice fails the additionality test if its adoption rate in the target region already exceeds 50\%. Krishi-Drishti evaluates this through a district-level database compiled from ICAR Annual Reports 2022--23 and NABARD District Agriculture Surveys.

\subsection{Cryptographic Traceability in Carbon Markets}
Double-counting — where the same carbon sequestration is credited multiple times — presents a major risk in voluntary carbon markets [14]. Cryptographic tools, like hash-chain ledgers and Merkle tree commitments, provide a lightweight tracking solution that does not rely on a live blockchain node. SHA-256 hash chaining means any alteration to a historical record invalidates the hash of every subsequent entry. Daily Merkle Root Anchors commit all events of the day into a single 64-character digest, which can then be published or anchored to a public Layer-2 blockchain (e.g., Polygon, Base) for external audits.

\subsection{Deep Learning for Crop Disease Detection}
While crop disease detection operates as a supporting feature, it helps sustain farmer engagement. Recent studies using CNNs for crop disease [7, 8, 9] report high accuracy but are often restricted to specific crops. Our approach pairs a lightweight MobileNetV2 classifier with a zero-shot Gemini LLM overlay to provide multi-crop capabilities and recommend localized, organic mitigation strategies.

% ═══════════════════════════════════════════════════════════════════════════════
\section{MATERIALS AND METHODS}
% ═══════════════════════════════════════════════════════════════════════════════

This section details the satellite data pipeline, the SOC ensemble model, the carbon credit calculation method, the additionality scoring system, the cryptographic token structure, and the supporting disease classification model.

\subsection{Sentinel-2 Satellite Data Pipeline and Spectral Feature Engineering}
The data source for SOC prediction is Sentinel-2 Level-2A (Bottom-of-Atmosphere) multispectral imagery, retrieved via the Google Earth Engine (GEE) Python API. For each farm plot, the pipeline applies a cloud-masking filter (QA60 band) and builds a median temporal composite over the previous 90-day period. Pixels with an NDVI $< 0.3$ are isolated for SOC analysis to ensure the model evaluates bare or sparsely vegetated soil [6].

The extracted spectral features include the Normalized Difference Vegetation Index (NDVI) as a mask, the Enhanced Vegetation Index (EVI) for atmospheric correction, the Bare Soil Index (BSI) for soil organic matter capture in the SWIR range, and the Modified Soil Adjusted Vegetation Index (MSAVI).

The final feature vector provided to the SOC model contains: NDVI, EVI, BSI, MSAVI, and the raw reflectance of Sentinel-2 Bands B2 (Blue), B4 (Red), B5, B6, B7 (Red-Edge), B8 (NIR), B11, B12 (SWIR). These 12 features are calculated for each plot as spatial medians over the valid bare-soil pixels.

\begin{figure}[H]
    \centering
    \fbox{\includegraphics[width=0.95\textwidth]{fig_satellite_pipeline.jpg}}
    \caption{Sentinel-2 Satellite SOC Prediction Pipeline for Carbon Credit Verification. The six-stage pipeline proceeds from raw L2A imagery ingestion through GEE cloud processing, spectral feature engineering, Stacking Ensemble ML (R\textsuperscript{2}=0.89), SOC stock estimation, and finally VM0042-compliant carbon credit calculation. \newline \centering Source: Authors, (2026).}
\end{figure}

\subsection{SOC Ensemble Model for Carbon Credit Quantification}
The SOC regression model relies on a Stacking Ensemble architecture, an approach shown to perform well for Sentinel-2 SOC mapping [15]. Three base learners are independently trained on the spectral features:
\begin{enumerate}
    \item \textbf{Random Forest (RF)}: 200 decision trees, max depth = 15. RF helps handle outliers and captures non-linear relationships.
    \item \textbf{XGBoost}: Gradient-boosted trees with a learning rate = 0.05, 300 estimators, and colsample\_bytree = 0.8.
    \item \textbf{CatBoost}: Gradient boosting optimized for categorical and mixed features, which is useful for processing varied soil-type contexts across districts.
\end{enumerate}
A Ridge Regression meta-learner is trained on the cross-validated out-of-fold predictions from these base models, producing the final SOC estimate in g/kg. The model was trained using 50 geo-referenced, laboratory-verified soil samples collected across Karnataka, Maharashtra, and Tamil Nadu, covering SOC ranges of 3.2--28.6 g/kg.

\subsection{VM0042 Carbon Credit Calculation Engine}
The engine converts satellite-predicted SOC values into carbon credits following the Verra VM0042 IALM methodology. It first calculates the SOC stock using the IPCC Tier 2 formula (incorporating SOC\% and Bulk Density). The net SOC change from baseline is multiplied by the farm area to estimate total sequestered carbon ($tC$), which is then converted to $tCO_2e$. We apply the Verra MRV net emission reduction logic ($ER_y = BE_y - PE_y - LE_y$), which includes a 10\% leakage deduction ($LE_y$). Finally, a 15\% buffer pool deduction is factored in to account for permanence risks, yielding the available credits.

\subsection{Four-Pillar MRV Evidence Framework}
The Carbon Vault employs a Four-Pillar readiness gate. Each pillar must be fulfilled before credits are issued. A \texttt{readiness} API endpoint checks each pillar and provides a structured audit report.

\textbf{Pillar 1 — Land \& Activity Proof (Verra VM0042 \S5 / Gold Standard LUF AGR FM):}
The farmer maps the GPS polygon boundary of the farm. The system checks if the regenerative practice is \textit{additional} by referencing a district-level database:
\begin{equation}
    \text{Additionality} = \begin{cases} \text{Eligible} & \text{if } \text{AdoptionRate}_{district} < 0.50 \\ \text{Rejected} & \text{otherwise} \end{cases}
\end{equation}
The farmer also maintains a Management Practice Log, detailing field activities (planting, fertilizer application, irrigation, harvest) with GPS coordinates and optional geotagged photos.

\textbf{Pillar 2 — Physical Measurement \& Sampling (IPCC M\&R Protocol):}
At least one GPS-tagged physical soil sample ($\geq$ 30 cm depth) is needed, alongside a laboratory accreditation certificate. The administrator verifies the submitted SOC\% and bulk density values. For projects older than three years, an additional Monitoring \& Re-measurement (M\&R) soil sample is required.

\textbf{Pillar 3 — Predictive Modeling \& Remote Sensing (Sentinel-2 dMRV):}
The Sentinel-2 GEE pipeline must have completed a satellite scan within the past 180 days, producing valid NDVI, BSI, and SOC estimates. A 10\% leakage deduction and a 15\% buffer pool deduction are automatically applied to this satellite-derived estimate.

\textbf{Pillar 4 — Third-Party Verification (CCTS/VVB equivalent):}
The farmer must provide at least one geotagged field photograph. Krishi-Drishti generates a formal Verification Report (\texttt{KD-VR-\{project\_id\}-\{date\}}) covering all four pillars. An administrator then approves or rejects the report, creating a documented \texttt{AdminCreditDecision} record for CCTS/BEE submission.

\begin{figure}[H]
    \centering
    \fbox{\includegraphics[width=0.95\textwidth]{fig_four_pillar_mrv.jpg}}
    \caption{Four-Pillar MRV Evidence Framework of the Krishi-Drishti Carbon Vault. Each pillar must independently pass its evidence gate before carbon credits can be issued. The framework is computationally enforced via a dedicated readiness API, aligned with Verra VM0042, IPCC M\&R Protocol, and CCTS/BEE standards. \newline \centering Source: Authors, (2026).}
\end{figure}

\subsection{HarvestToken: Cryptographic Credit Provenance}
To address double-counting concerns and document credit ownership, approved credits are minted as a \texttt{HarvestToken} (ID format: \texttt{KD-HTK-YYYY-NNNNN}). The token uses a SHA-256 digest to chain agronomic details with the preceding token hash:
\begin{equation}
    H_n = SHA256\,(crop\_type_n \| yield_n \| GPS_n \| NDVI_n \| CO_{2eq,n} \| H_{n-1})
\end{equation}
where $H_{n-1}$ is the hash of the preceding token. Altering a historical token will break the chain, making tampering detectable. A daily \textbf{Merkle Root Anchor} consolidates all daily crop cycle events into a 64-character SHA-256 root digest, allowing future integration with public Layer-2 blockchains.

\begin{figure}[H]
    \centering
    \fbox{\includegraphics[width=0.95\textwidth]{fig_harvest_token.jpg}}
    \caption{HarvestToken SHA-256 Hash Chain Architecture. Each token encapsulates crop provenance, GPS coordinates, Sentinel-2 NDVI, CO\textsubscript{2}e carbon data, and the hash of the preceding token --- forming a tamper-evident linked ledger. Any modification to block $H_n$ immediately invalidates all subsequent blocks $H_{n+1} \ldots H_N$. A Daily Merkle Root Anchor provides an external blockchain anchoring pathway. \newline \centering Source: Authors, (2026).}
\end{figure}

The token records crop variety, harvest GPS coordinates, yield, harvested area, carbon footprint, harvest-time NDVI, chemical input logs, associated carbon credits, and farming methodology. This data format supports external auditing without manual document reviews.

\subsection{District-Level Additionality Scoring Engine}
The additionality database relies on the ICAR Annual Report 2022--23, NABARD District Agriculture Survey Reports, and Ministry of Agriculture \& Farmers' Welfare State Statistics 2023. For each district $d$ and regenerative practice $p$, a regional adoption rate $A(d,p) \in [0,1]$ is tracked. The evaluation follows Verra VM0042 and CCTS/BEE guidelines:
\begin{equation}
    \text{isAdditional}(d,p) = \mathbf{1}\left[A(d,p) < 0.50\right]
\end{equation}
Practices evaluated include No-Till, Cover-Crop, Agroforestry, Composting, Biochar, Green-Manure, Reduced-Tillage, Mulching, Crop-Rotation, and Irrigation-Optimization.

\subsection{Supporting Module: MobileNetV2 Crop Disease Classifier}
A MobileNetV2 classifier, trained on a 4,000-image subset of the PlantVillage dataset across 5 classes, serves as the local inference model for disease diagnosis. The classification output is passed to a Google Gemini 2.5 Flash LLM, which generates localized and organic-first treatment recommendations.

% ═══════════════════════════════════════════════════════════════════════════════
\section{PROPOSED SYSTEM ARCHITECTURE OF KRISHI-DRISHTI}
% ═══════════════════════════════════════════════════════════════════════════════

Krishi-Drishti is constructed with a React/TypeScript frontend and a FastAPI Python backend. The architecture centers around the Carbon Vault, with other components feeding data and engagement into the verification pipeline.

\subsection{Carbon Vault: End-to-End Credit Lifecycle}
The Carbon Vault oversees the credit lifecycle: Enrollment $\rightarrow$ Evidence Collection $\rightarrow$ Satellite Monitoring $\rightarrow$ Readiness Gate $\rightarrow$ Admin Verification $\rightarrow$ Credit Issuance $\rightarrow$ Wallet Claim $\rightarrow$ Payout.

\begin{figure}[H]
    \centering
    \fbox{\includegraphics[width=0.95\textwidth]{fig_carbon_lifecycle.jpg}}
    \caption{Carbon Credit Issuance Lifecycle of the Krishi-Drishti Carbon Vault. The 8-stage pipeline proceeds from Farm Enrollment and Additionality Checking, through Satellite Monitoring and the 4-Pillar MRV Gate, to auto-generated Verification Report, Admin Review, HarvestToken Minting, and final Farmer Wallet Payout at \rupee{}1,200/tCO\textsubscript{2}e (80\% farmer share). \newline \centering Source: Authors, (2026).}
\end{figure}

Upon enrollment, the farmer specifies a methodology and a registered plot. The system checks additionality and retrieves satellite data. If verified, the project status becomes \texttt{Enrolled}. The farmer then fulfills the Four Pillars using the app—logging activities, uploading certificates, and providing photos. A readiness dashboard tracks progress.

Once all pillars pass, a Verification Report is generated containing the GPS boundary, additionality score, Sentinel-2 SOC estimate, soil sample results, photo evidence, and a SHA-256 integrity hash. An administrator evaluates the report to approve or reject the issuance.

\subsection{Carbon Wallet and Aggregator Integration}
Approved credits are deposited into the Carbon Wallet, which displays verified credits (tCO$_2$e), available credits, locked credits (buffer pool), and estimated value in INR (at \rupee{}1,200/tCO$_2$e). KYC-verified farmers can request payouts for 80\% of the gross value, with the remaining 20\% allocated as a platform fee.

The platform connects to four carbon aggregators:
\begin{itemize}
    \item \textbf{Krishi-Drishti Platform (GCP Beta)}: 15\% fee, 14-day settlement, India Green Credit Program.
    \item \textbf{TrayamBhu Tech (CCTS/BEE)}: 18\% fee, 21-day settlement, domestic compliance market.
    \item \textbf{Grow Indigo (Verra VCS/Gold Standard)}: 20\% fee, 30-day settlement, international voluntary market.
    \item \textbf{Boomitra (Verra VCS)}: 25\% fee, 45-day settlement, premium international buyers.
\end{itemize}

\subsection{Supporting AI Vision and Disease Diagnostic Pipeline}
The diagnostic pipeline involves three steps: (1) Image Acquisition via smartphone; (2) Local Inference using MobileNetV2; and (3) Generative Remediation, where the Gemini 2.5 Flash LLM provides localized treatment plans, prioritizing organic solutions for farmers in the carbon program.

\subsection{Digital Twin for Carbon Visualization}
A WebGL-based Digital Twin offers a 3D farm model layered with false-color NDVI heatmaps, topographical moisture data, and SOC depth-maps. This tool helps visualize the satellite verification process and provides additional context during the Pillar 4 admin review.

% ═══════════════════════════════════════════════════════════════════════════════
\section{RESULTS AND DISCUSSIONS}
% ═══════════════════════════════════════════════════════════════════════════════

\subsection{Satellite SOC Model Performance}
The Stacking Ensemble SOC model was evaluated using 5-fold cross-validation on a held-out test set from the 50 sampled plots. Table 2 outlines the performance of the base models and the ensemble.

\begin{table}[H]
\centering
\caption{Performance Comparison of Satellite-Based SOC Prediction Models on 50 Indian Farmland Plots.}
\begin{tabular}{|l|c|c|c|}
\hline
\textbf{Model} & \textbf{R$^2$ Score} & \textbf{RMSE (g/kg)} & \textbf{MAE (g/kg)} \\
\hline
Random Forest (Standalone) & 0.82 & 1.94 & 1.51 \\
XGBoost (Standalone) & 0.85 & 1.79 & 1.38 \\
CatBoost (Standalone) & 0.83 & 1.88 & 1.44 \\
\textbf{Stacking Ensemble (RF + XGB + CAT)} & \textbf{0.89} & \textbf{1.52} & \textbf{1.19} \\
\hline
\end{tabular}
\end{table}

The ensemble's R$^2$ of 0.89 indicates that a large portion of the variance in physical SOC measurements can be explained by the satellite spectral features. The RMSE of 1.52 g/kg falls within the tolerance margins typically required for the Verra voluntary carbon market ($\leq$ 3 g/kg). This result suggests that Sentinel-2 imagery can serve as a primary basis for carbon credit calculations in this context. Feature importance metrics showed that BSI, SWIR Band 11 (B11), and Red-Edge Band 5 (B5) provided the most predictive value, while NDVI primarily served as a pre-filtering mask.

\subsection{Carbon Vault MRV Framework Validation}
The Four-Pillar MRV system was tested with 12 simulated farm enrollments involving various methodologies and plot sizes (0.5--3.2 hectares) in 6 districts. Table 3 presents the validation results.

\begin{table}[H]
\centering
\caption{Four-Pillar MRV Validation Results Across 12 Farm Enrollments.}
\begin{tabular}{|l|c|c|}
\hline
\textbf{MRV Component} & \textbf{Result} & \textbf{Standard} \\
\hline
Additionality Classification Accuracy & 91.7\% (11/12) & Verra VM0042 / CCTS \\
Pillar 1 Pass Rate (GPS + Activity Log) & 100\% (12/12) & VM0042 \S5 \\
Pillar 2 Pass Rate (Soil Sample) & 83.3\% (10/12) & IPCC M\&R \\
Pillar 3 Pass Rate (Satellite Scan $<$ 180d) & 100\% (12/12) & Sentinel-2 dMRV \\
Pillar 4 Pass Rate (Evidence + Admin Review) & 91.7\% (11/12) & CCTS/VVB \\
End-to-End Verification Latency & $<$ 48 hours & vs. 6--18 months (traditional) \\
HarvestToken Hash Mismatches & 0 / 47 tokens & SHA-256 integrity \\
\hline
\end{tabular}
\end{table}

The one additionality rejection correctly flagged a No-Till project in Nashik district, where the adoption rate was 0.62 (above the 0.50 threshold). The end-to-end verification time averaged under 48 hours. Furthermore, all 47 minted HarvestTokens retained their hash integrity.

\subsection{Carbon Credit Yield and Economic Impact}
Table 4 estimates the satellite-verified carbon credit yield and financial impact for four supported regenerative methodologies.

\begin{table}[H]
\centering
\caption{Carbon Vault: Satellite-Verified Credit Yield and Farmer Income Impact by Methodology.}
\begin{tabular}{|l|c|c|c|c|c|}
\hline
\textbf{Methodology} & \textbf{Plot (ha)} & \textbf{Gross tCO$_2$e/yr} & \textbf{After 15\% Buffer} & \textbf{INR Payout} & \textbf{USD Equiv.} \\
\hline
No-Till & 1.1 & 4.10 & 3.49 & \rupee{}3,349 & \$40 \\
Cover-Crop & 1.1 & 5.30 & 4.51 & \rupee{}5,412 & \$65 \\
Agroforestry & 2.0 & 8.40 & 7.14 & \rupee{}8,568 & \$102 \\
Composting & 0.8 & 2.80 & 2.38 & \rupee{}2,856 & \$34 \\
\hline
\textbf{Average (1.1 ha)} & \textbf{1.1} & \textbf{4.41} & \textbf{3.75} & \textbf{\rupee{}4,500} & \textbf{\$54} \\
\hline
\end{tabular}
\end{table}

For an average farm size of 1.1 hectares, the system outlines a potential supplemental income of USD \$52--\$132 per year. This income relies on verifiable sustainable practices rather than significant new capital expenditures.

\subsection{Additionality Engine Validation}
The additionality engine was evaluated against 12 district-practice scenarios, classifying 11 of 12 correctly (91.7\% accuracy). The exception was a borderline composting case in Pune district (adoption rate = 0.48), which was classified as additional but considered ambiguous by agronomists.

\subsection{AI Vision Diagnosis Performance (Supporting Module)}
The MobileNetV2 classifier achieved an 88.1\% overall accuracy on 730 test images. This capability supports the Gemini LLM advisory function, offering farmers relevant treatment recommendations while highlighting organic options for those in carbon projects.

\subsection{Comparative Analysis}
Table 6 compares Krishi-Drishti with similar systems documented in recent literature.

\begin{table}[H]
\centering
\caption{Comparison of Krishi-Drishti Against Related Agricultural AI Systems in Literature.}
\begin{tabularx}{\textwidth}{|p{2.2cm}|X|c|p{3.5cm}|}
\hline
\textbf{System / Reference} & \textbf{Core Technology} & \textbf{Accuracy} & \textbf{Limitation vs. Krishi-Drishti} \\
\hline
Iftikhar et al. [7] & Enhanced CNN, 3 crops & 98.1\% & Disease-only; no SOC or carbon credit. \\
\hline
Hasan et al. [8] & Lightweight CNN, Rice & 97.9\% & Single crop; no market or carbon linkage. \\
\hline
Saha et al. [21] & IoT + RF Precision Ag & 97.3\% & Requires physical IoT hardware; no carbon MRV. \\
\hline
Iswarya \& Tamizhazhagan [12] & QRF + Sentinel-2 SOC & High R$^2$ & SOC prediction only; no credit issuance pipeline. \\
\hline
\textbf{Krishi-Drishti (Ours)} & \textbf{4-Pillar dMRV + Sentinel-2 SOC + SHA-256 HarvestToken + VM0042} & \textbf{R$^2$=0.89 (SOC); 88.1\% (Vision)} & \textbf{Full end-to-end carbon credit verification. No physical hardware required.} \\
\hline
\end{tabularx}
\end{table}

A key feature of Krishi-Drishti is its ability to support the complete carbon credit lifecycle—from satellite SOC estimation to credit minting—in one application designed around standard MRV practices.

% ═══════════════════════════════════════════════════════════════════════════════
\section{CONCLUSIONS}
% ═══════════════════════════════════════════════════════════════════════════════

This paper detailed Krishi-Drishti, a mobile-first platform focused on lowering the barriers to agricultural carbon credit verification through a satellite-driven Digital MRV system. Our results indicate that Sentinel-2 multispectral imagery, when evaluated using a Stacking Ensemble of Random Forest, XGBoost, and CatBoost models via Google Earth Engine, can produce reliable Soil Organic Carbon estimates (R$^2$ = 0.89, RMSE = 1.52 g/kg). These metrics align with the evidence standards for Verra VM0042 IALM and India's CCTS/BEE frameworks while reducing the reliance on physical soil testing.

Validation of the Four-Pillar MRV framework across 12 farm enrollments showed that the entire verification process, from GPS registration to cryptographic HarvestToken minting, can be completed in under 48 hours. This significantly shortens the typical timeline of traditional audits and bypasses the high laboratory costs that previously prevented smallholder farmers from accessing the carbon market.

Our tests on the additionality engine (91.7\% accuracy) and the SHA-256 HarvestToken hash chain (no integrity violations observed) suggest that an end-to-end issuance system is both workable and secure. Economically, this model could provide farmers with an estimated USD \$52--\$132 annually. The integration of extra tools, like MobileNetV2 for disease diagnosis and a peer-to-peer marketplace, also encourages continuous engagement with sustainable practices.

Future research will aim to expand the SOC model to cover more agro-climatic zones in India, incorporate Landsat 8/9 data for better temporal resolution, and anchor Merkle Root digests to public Layer-2 blockchains for improved transparency. We also plan a broader field trial involving at least 500 farms across five states to further evaluate the platform's scalability.

\section{AUTHOR'S CONTRIBUTION}
\textbf{Dr. Rama Krishna K:} Supervision, Project Administration. \\
\textbf{Dr. Sheethal Aji Mani:} Methodology, Validation, Writing -- Review and Editing. \\
\textbf{Bhagesh Biradar:} Conceptualization, Carbon Vault Architecture, Software Development, Methodology, Writing -- Original Draft. \\
\textbf{Manjunath Khot:} Data Curation, Satellite Pipeline Development, Formal Analysis. \\
\textbf{Jayanth R:} Investigation, Field Testing, Resources. \\
\textbf{Pruthviraja S:} Visualization, Digital Twin Development, UI/UX Design.

\section{ACKNOWLEDGMENTS}
The authors acknowledge the use of Google Earth Engine for satellite data processing, the Sentinel-2 Level-2A data made freely available by the European Space Agency (ESA) Copernicus Programme, the Google Gemini API for generative AI capabilities, and the Open-Meteo API for weather data. We extend our gratitude to the Visvesvaraya Technological University for supporting this project, and to the farmers of Karnataka and Maharashtra who participated in the soil sampling validation exercise.

\section{REFERENCES}
\noindent [1] Bhatia, A. (2025). ``An AI-Powered Agricultural Decision Support System for Smallholder Farmers in India.'' \textit{Smart Agricultural Technology}, SSRN-6944006.

\noindent [2] Thilakarathne, N.N., et al. (2022). ``A Cloud Enabled Crop Recommendation Platform for Machine Learning-Driven Precision Farming.'' \textit{Sensors}, 22(16), 6299.

\noindent [3] Gokul, S., et al. (2026). ``Carbon credit mechanisms for sustainable agriculture and opportunities in North East India.'' \textit{Discover Sustainability}, 7, 857.

\noindent [4] Abiri, R., et al. (2023). ``Application of digital technologies for ensuring agricultural productivity.'' \textit{Heliyon}, 9, e22601.

\noindent [5] Nawaz, M., \& Babar, M.I.K. (2025). ``IoT and AI for smart agriculture in resource-constrained environments.'' \textit{Discover Internet of Things}, 5, 24.

\noindent [6] Vaudour, E., et al. (2022). ``Satellite Imagery to Map Topsoil Organic Carbon Content over Cultivated Areas: An Overview.'' \textit{Remote Sensing}, 14, 2917.

\noindent [7] Iftikhar, M., et al. (2024). ``Plant disease management: a fine-tuned enhanced CNN approach with mobile app integration.'' \textit{Artificial Intelligence Review}, 57, 167.

\noindent [8] Hasan, M.M., et al. (2023). ``Enhancing Rice Crop Management: Disease Classification Using CNNs and Mobile Application Integration.'' \textit{Agriculture}, 13(8), 1549.

\noindent [9] Askale, G.T., et al. (2025). ``Mobile based deep CNN model for maize leaf disease detection and classification.'' \textit{Plant Methods}, 21, 72.

\noindent [10] Khaddi, E., et al. (2026). ``Shaping smarter prediction: A literature review of GIS, remote sensing and ML applications in digital SOC mapping.'' \textit{Turkish Journal of Remote Sensing}, 8, 1808033.

\noindent [11] Xiao, X., et al. (2024). ``Environmental variables improve the accuracy of remote sensing estimation of soil organic carbon content.'' \textit{Scientific Reports}, 14, 18964.

\noindent [12] Iswarya, R., \& Tamizhazhagan, V. (2026). ``Satellite-Derived Prediction of Soil Organic Carbon Under Bare-Soil Conditions Using Sentinel-2 Spectral Indices and Quantile Regression Forests.'' \textit{International Journal of AI and Machine Learning}, 6(6s), 76.

\noindent [13] Vaudour, E., et al. (2022). ``Satellite Imagery to Map Topsoil Organic Carbon Content over Cultivated Areas: An Overview.'' \textit{Remote Sensing}, 14, 2917.

\noindent [14] Hamrick, K. (2023). ``State of Voluntary Carbon Markets 2023.'' \textit{Ecosystem Marketplace}, Forest Trends.

\noindent [15] Triantakonstantis, D., \& Karakostas, A. (2025). ``Soil Organic Carbon Monitoring and Modelling via Machine Learning Methods Using Soil and Remote Sensing Data.'' \textit{Agriculture}, 15, 910.

\noindent [16] Kershenbaum, A., et al. (2024). ``Automatic detection for bioacoustic research: a practical guide from and for biologists and computer scientists.'' \textit{Biological Reviews}. DOI: 10.1111/brv.13155.

\noindent [17] Hearon, L.E., et al. (2025). ``buzzdetect: an open-source deep learning tool for automated bioacoustic pollinator monitoring.'' \textit{Journal of Insect Science}, 25(6), ieaf104.

\noindent [21] Saha, G., et al. (2025). ``Smart IoT-driven precision agriculture: Land mapping, crop prediction, and irrigation system.'' \textit{PLOS ONE}, 20(3), e0319268.

\end{document}
